\documentclass[10pt,landscape]{article}

\usepackage[letterpaper,margin=0.18in]{geometry}
\usepackage{amsmath,amssymb,mathtools,array}
\usepackage[T1]{fontenc}
\usepackage[utf8]{inputenc}
\usepackage{lmodern}
\usepackage{multicol}
\usepackage{enumitem}

\pagestyle{empty}
\setlength{\parindent}{0pt}
\setlength{\parskip}{0pt}
\setlength{\columnsep}{7pt}
\setlength{\premulticols}{2pt}
\setlength{\postmulticols}{2pt}
\setlength{\multicolsep}{2pt}
\setlength{\abovedisplayskip}{2pt}
\setlength{\belowdisplayskip}{2pt}
\setlength{\abovedisplayshortskip}{1pt}
\setlength{\belowdisplayshortskip}{1pt}
\setlist[itemize]{leftmargin=*, itemsep=0pt, topsep=1pt, parsep=0pt, partopsep=0pt}
\setlist[enumerate]{leftmargin=*, itemsep=0pt, topsep=1pt, parsep=0pt, partopsep=0pt}

\makeatletter
\renewcommand{\section}{\@startsection{section}{1}{0mm}%
  {-0.6ex plus -.2ex minus -.1ex}%
  {0.4ex plus .1ex}%
  {\normalfont\normalsize\bfseries}}
\renewcommand{\subsection}{\@startsection{subsection}{2}{0mm}%
  {-0.5ex plus -.2ex minus -.1ex}%
  {0.2ex plus .1ex}%
  {\normalfont\small\bfseries}}
\makeatother

\newcommand{\E}{\mathbb{E}}
\newcommand{\KL}{\operatorname{KL}}
\newcommand{\clip}{\operatorname{clip}}
\newcommand{\InfoNCE}{\operatorname{InfoNCE}}
\newcommand{\head}[1]{\section*{#1}}
\newcommand{\mini}[1]{\textbf{#1}}
\newcommand{\defs}[1]{{\emph{#1}}\par}
\newcommand{\intuition}[1]{{\scriptsize\emph{Intuition: #1}\par}}

\begin{document}
\raggedright
\footnotesize

\begin{center}
\textbf{6.8610 Quiz 2 Cheatsheet}
\end{center}

\mini{Question cues:} bottleneck first; watch absolute words (\emph{always}, \emph{only}, \emph{guarantees}, \emph{equivalent}); reject answers by missing property; say quality \emph{and} systems constraint; simplify toy objectives before interpreting them.

\begin{multicols}{3}

\head{Post-Training Flow}
\[
\text{pretrain}
\rightarrow
\text{mid-train}
\rightarrow
\text{SFT}
\rightarrow
\text{prefs}
\rightarrow
\text{RL}
\]
\intuition{Move from broad language knowledge to sharper task behavior and reward optimization.}

\mini{Mid-train}
\begin{itemize}
\item Continue next-token prediction on smaller, higher-value data: code, books, technical text, long-context/tool traces.
\item Shorter ctx can be cheaper early; expand ctx later once basic priors exist.
\item Later data matters more: better to forget early boilerplate than late high-value data.
\end{itemize}

\mini{SFT objective}
\defs{\(\theta=\) params,\; \(\mathcal D=\{(x,y)\}\) demo pairs,\; \(x=\) prompt,\; \(y=\) assistant response.}
\[
\max_\theta \sum_{(x,y)\in\mathcal D}\log p_\theta(y\mid x)
\]
\intuition{Increase probability of desired assistant completions.}
\begin{itemize}
\item Usually mask loss on prompt/user tokens; supervise mainly assistant tokens.
\item Teaches instruction following, chat format, refusals, response structure.
\item Extracts latent skills; does not replace pretraining.
\end{itemize}

\mini{SFT failure modes}
\begin{itemize}
\item Label bottleneck.
\item One gold answer when many would work.
\item Can teach style of competence more than calibrated truth.
\item ``I don't know'' is hard from demos alone.
\end{itemize}

\head{Preference Learning}
\mini{Setup}
\defs{\(x=\) prompt,\; \(y_w/y_\ell=\) winner/loser,\; \(p_\theta=\) policy,\; \(p_0=\) reference,\; \(\sigma=\) logistic.}

\mini{Bradley--Terry}
\defs{\(r_\theta(x,y)=\) learned scalar preference score.}
\[
P(y_w\succ y_\ell\mid x)=\sigma(r(x,y_w)-r(x,y_\ell))
\]
\intuition{A winner is more likely when its reward exceeds the loser's.}
\[
\mathcal L_{\mathrm{BT}}
=
-\log \sigma(r_\theta(x,y_w)-r_\theta(x,y_\ell))
\]
\intuition{Train the reward model to rank winners above losers.}

\mini{DPO shorthand}
\[
\Delta_\theta
=
\log\frac{p_\theta(y_w\mid x)}{p_0(y_w\mid x)}
-
\log\frac{p_\theta(y_\ell\mid x)}{p_0(y_\ell\mid x)}
\]
\intuition{Measure how much more the policy prefers the winner than the loser relative to the reference.}
\[
\mathcal L_{\mathrm{DPO}}=-\log \sigma(\lambda\Delta_\theta)
\]
\intuition{Push the policy toward preferred answers while keeping it reference-anchored.}
\begin{itemize}
\item Uses LM probs relative to reference; no separate reward model required.
\item Reference model limits drift while satisfying winner/loser constraints.
\item Pairwise labels are often easier than writing ideal completions.
\end{itemize}

\mini{Reward parameterization behind DPO}
\defs{\(\lambda=\) preference scale,\; \(Z(x)=\) prompt-wise partition function.}
\[
r_\theta(x,y)
=
\lambda\log\frac{p_\theta(y\mid x)}{p_0(y\mid x)}
+ \lambda\log Z(x)
\]
\intuition{Treat reward as scaled log-prob improvement over the reference model.}

\mini{If SLiC fits in memory}
\defs{Here \(\Delta>0\) is a margin, not the DPO \(\Delta_\theta\) above.}
\[
\max(0,\Delta+\log p_\theta(y_\ell\mid x)-\log p_\theta(y_w\mid x))
\]
\intuition{Enforce a margin so the winner's log-prob beats the loser's.}

\mini{Keep straight}
\begin{itemize}
\item SFT = demonstrations.
\item DPO = preferences + reference.
\item RL = optimize reward on fresh model samples.
\end{itemize}

\head{RLHF Basics}
\mini{Reward-regularized objective}
\defs{\(x=\) prompt,\; \(y=\) completion,\; \(p_0=\) reference policy,\; \(\lambda=\) KL weight.}
\[
J(\theta)
=
\E_x\E_{y\sim p_\theta(\cdot\mid x)}[r(x,y)]
-\lambda \KL(p_\theta\|p_0)
\]
\intuition{Maximize reward, but penalize drifting too far from the reference policy.}

\mini{Reward model}
\defs{\(\phi=\) reward-model params,\; \(h_\phi(x,y)=\) features,\; \(w=\) linear head.}
\[
r_\phi(x,y)=w^\top h_\phi(x,y)
\]
\intuition{Map each completion to a single learned scalar score.}
\[
\max_\phi \log \sigma(r_\phi(x,y_w)-r_\phi(x,y_\ell))
\]
\intuition{Fit that score so preferred completions outrank dispreferred ones.}

\mini{Vanilla policy gradient}
\defs{\(r(x,y)=\) scalar reward.}
\[
\nabla_\theta \E_{y\sim p_\theta}[r]
=
\E_{y\sim p_\theta}
\left[\nabla_\theta\log p_\theta(y\mid x)\, r(x,y)\right]
\]
\intuition{Increase probability of sampled actions that earned high reward.}

\mini{Value / advantage}
\[
V(x)=\E_{y\sim p_\theta}[r(x,y)],
\qquad
A(x,y)=r(x,y)-V(x)
\]
\intuition{Advantage tells whether a sampled completion was better than expected.}
\begin{itemize}
\item Baseline reduces variance; does \emph{not} change expected gradient.
\item Vanilla PG is unbiased, high variance, on-policy.
\item Big RL story: learn reward from scarce labels, then use compute on new samples.
\end{itemize}

\mini{KL-folding approx.}
\defs{\(r'(x,y)=\) reward after folding the reference-model KL term into the scalar reward.}
\[
r'(x,y)=r(x,y)-\log p_\theta(y\mid x)+\log p_0(y\mid x)
\]
\intuition{Absorb the KL term into a modified per-sample reward.}

\mini{Tiny cue}
\begin{itemize}
\item Equal SFT probs, rewards \(1/0\), \(\beta=\) KL weight \(=1\) \(\Rightarrow\) better response gets prob \(e/(e+1)\).
\end{itemize}

\head{PPO / GRPO}
\mini{Importance ratio}
\defs{\(p_{\theta_{\mathrm{old}}}=\) old policy,\; \(A=\) advantage,\; \(\epsilon=\) clip radius.}
\[
\rho(x,y)=\frac{p_\theta(y\mid x)}{p_{\theta_{\mathrm{old}}}(y\mid x)},
\qquad
\bar\rho=\clip(\rho,1-\epsilon,1+\epsilon)
\]
\intuition{Compare new and old policies, then cap oversized probability changes.}

\mini{PPO objective}
\[
\mathcal L_{\mathrm{PPO}}
=
\E_{y\sim p_{\theta_{\mathrm{old}}}}
\left[
\min\left(\rho A,\bar\rho A\right)
\right]
\]
\intuition{Improve good actions without letting one update move the policy too far.}
\begin{itemize}
\item Sample reuse via importance sampling.
\item Advantage weighting.
\item Clipping stabilizes policy updates.
\end{itemize}

\mini{GRPO}
\defs{\(r_i=\) rewards for samples \(i\) from one prompt,\; \(\mu,\sigma=\) group mean/std.}
\[
\mu=\operatorname{mean}(r_i),
\qquad
\sigma=\operatorname{std}(r_i),
\qquad
A_i=\frac{r_i-\mu}{\sigma}
\]
\intuition{Turn raw rewards into within-prompt relative advantages.}
\begin{itemize}
\item Same PPO-style clipping, but no learned value model.
\item Group normalization is convenient when many samples come from one prompt.
\end{itemize}

\mini{Fast compare}
\begin{itemize}
\item \textbf{SFT:} imitate demos.
\item \textbf{DPO:} fit preferences against reference.
\item \textbf{PPO:} practical RL workhorse.
\item \textbf{GRPO:} simpler RL variant with within-group normalization.
\end{itemize}

\mini{What usually matters}
\begin{itemize}
\item supervision quality
\item reward credit
\item stability
\item whether fresh model samples are needed
\end{itemize}

\end{multicols}

\newpage

\begin{center}
\textbf{6.8610 Quiz 2 Cheatsheet}
\end{center}

\mini{Retriever tradeoffs}\par
\setlength{\tabcolsep}{4pt}
\renewcommand{\arraystretch}{1.05}
\begin{tabular*}{\textwidth}{@{\extracolsep{\fill}} l l l l l @{}}
\textbf{Method} & \textbf{Interaction} & \textbf{Precompute} & \textbf{Query cost} & \textbf{Best use} \\
XE & full \(q\)-\(d\) & none & \(O(N)\) TF calls & rerank \\
DE & 1 vec/doc & docs offline & \(O(1)\) TF + \(O(N)\) dots & first stage \\
Late int. & token-level & doc toks offline & small q TF + ANN & middle ground \\
ANN / IVF & prune candidates & index offline & sublinear-ish shortlist & scale \\
\end{tabular*}

\begin{multicols}{3}

\head{Neural Retrieval}
\mini{Dual-encoder score}
\defs{\(q=\) query,\; \(d=\) doc,\; \(d^+=\) positive doc,\; \(d_i^-=\) negative docs.}
\[
s(q,d)=f(q)^\top g(d)
\]
\intuition{Retrieve by embedding similarity between the query and each document.}
\[
s^+=s(q,d^+),
\qquad
s_i^-=s(q,d_i^-)
\]
\intuition{Name the positive-match score and the negative distractor scores.}

\mini{InfoNCE}
\defs{\(\tau=\) temperature,\; \(s^+=\) positive score,\; \(s_i^-=\) negative scores.}
\[
\mathcal L_{\InfoNCE}
=
-\log
\frac{
\exp(s^+/\tau)
}{
\exp(s^+/\tau)
+
\sum_i \exp(s_i^-/\tau)
}
\]
\intuition{Make the positive document score beat many negatives in the softmax.}

\mini{Late interaction / ColBERT}
\defs{\(E_{q_i},E_{d_j}=\) query/doc token embeddings.}
\[
S(q,d)=\sum_i \max_j E_{q_i}^\top E_{d_j}
\]
\intuition{Each query token takes credit from its best-matching document token.}

\mini{Retrieval reminders}
\begin{itemize}
\item Cross-encoders are great rerankers, bad first-stage retrievers.
\item Hard negatives matter more than easy random negatives.
\item Single-vector dense retrieval often fails on OOD generalization.
\item Late interaction keeps precompute \emph{and} token-level matching.
\item ANN still works: high score needs at least one close token vector.
\end{itemize}

\head{RAG}
\mini{Closed-book} answer from params only.\par
\mini{Open-book} retrieve docs and answer from evidence.

\mini{RAG factorization}
\defs{\(q=\) query,\; \(a=\) answer,\; \(d=\) doc,\; \(\mathcal D=\) corpus,\; \(p_\eta=\) retriever,\; \(p_\theta=\) reader.}
\[
p(a\mid q)=\sum_{d\in\mathcal D} p_\eta(d\mid q)\, p_\theta(a\mid q,d)
\]
\intuition{Answer by averaging reader predictions over possible retrieved documents.}
\[
\text{practice: approximate over top-}K
\]

\mini{Reader choice}
\begin{itemize}
\item extractive = span; generative = decode over docs
\end{itemize}

\mini{Retriever supervision}
\begin{itemize}
\item Often only \((q,a)\), not gold docs; answer-containing doc is only a proxy.
\item Search relevance and answer usefulness are related, not identical.
\end{itemize}

\mini{Credit assignment}
\begin{itemize}
\item Right doc in top-\(K\) + wrong answer \(\Rightarrow\) reader/evidence-use failure.
\item Latent doc marginalization separates retriever success from reader failure; junk top-\(K\) can still push memorization.
\end{itemize}

\mini{FiD}
\begin{itemize}
\item encode passages separately; decoder attends across all
\end{itemize}

\head{Agents}
\mini{Inference scaling}
\begin{itemize}
\item Parallel: pass@\(k\), best-of-\(n\), self-consistency. Sequential: longer CoT / ``Wait'' forcing.
\end{itemize}

\mini{Pass@1 vs pass@\(k\)}
\begin{itemize}
\item pass@1 = one rollout succeeds.
\item pass@\(k\) = at least one of \(k\) rollouts succeeds.
\end{itemize}

\mini{Agent loop}
\defs{\(s_t=\) state,\; \(a_t=\) action,\; \(o_{t+1}=\) observation.}
\[
s_t \to a_t \to o_{t+1} \to s_{t+1}
\]
\intuition{Act, observe the environment, update state, and repeat.}

\mini{Tool use}
\begin{itemize}
\item emit tool call, execute externally, append observation, continue generation
\end{itemize}

\mini{ReAct vs CodeAct}
\begin{itemize}
\item \textbf{ReAct:} interleave reasoning and actions.
\item \textbf{CodeAct:} act by writing code; more general, more dangerous.
\end{itemize}

\mini{Long horizon}
\begin{itemize}
\item context rot + stateful environments
\item compaction keeps active context small while preserving task state
\item mention environment interaction and cost, not just answer quality
\end{itemize}

\head{Decoding}
\mini{Greedy trap}
\begin{itemize}
\item local argmax \(\neq\) best full sequence
\end{itemize}

\mini{Best-of-\(n\)}
\defs{\(p=\) base policy,\; \(q=\) best-of-\(n\) induced policy,\; \(n=\) number of samples.}
\[
\KL(q\|p)\le \log n
\]
\intuition{Best-of-\(n\) helps, but policy shift grows only logarithmically with \(n\).}
\begin{itemize}
\item sample \(n\) full outputs, keep best by reward/verifier; strong baseline
\item search error \(\neq\) model error; large policy shift needs exponential \(n\)
\end{itemize}

\head{Speculative / Contrastive}
\mini{Speculative}
\defs{\(p_t=\) target next-token dist,\; \(q_t=\) draft dist,\; \(\alpha=\) accept prob,\; \(A=\) accept rate.}
\[
\alpha(x)=\min\left(1,\frac{p_t(x)}{q_t(x)}\right),
\qquad
A=1-\tfrac{1}{2}\lVert p_t-q_t\rVert_1
\]
\intuition{Accept draft tokens when target agrees; speed rises as draft matches target.}
\begin{itemize}
\item draft \(q_t\) proposes, target \(p_t\) verifies; final token law still \(p_t\)
\item best when draft is close to target; bigger draft can raise acceptance but lose speed
\end{itemize}

\mini{Contrastive}
\defs{\(p_t^+=\) expert,\; \(p_t^-=\) amateur,\; \(\lambda=\) contrast penalty.}
\[
s_t(x)=\log p_t^+(x)-\lambda\log p_t^-(x)
\]
\intuition{Favor tokens the expert likes but the amateur does not over-prefer.}
\[
\max_x s_t(x)\equiv \max_x \frac{p_t^+(x)}{(p_t^-(x))^\lambda}
\]
\intuition{The same contrastive rule written as a ratio objective.}
\begin{itemize}
\item small \(\lambda\): near expert decoding
\item large \(\lambda\): penalize generic amateur-favored tokens; too large hurts coherence
\end{itemize}

\head{LM Programming}
\mini{Natural-language signatures}
\begin{itemize}
\item specify \emph{what}, not model-specific \emph{how}; e.g.\ \texttt{thread -> events, actions}
\end{itemize}

\mini{Three ingredients}
\begin{itemize}
\item code / modular structure
\item natural-language specs
\item data / feedback
\end{itemize}

\mini{Compile intent}
\begin{itemize}
\item prompts, weights (FT / RL), inference pipeline
\end{itemize}

\mini{Intervention map}
\begin{itemize}
\item underspecified task \(\Rightarrow\) modularize + prompt/program optimization
\item missing knowledge \(\Rightarrow\) retrieval, data, or weights
\item prompt optimization cannot add knowledge; it complements FT / RL
\end{itemize}

\head{Tokenization}
\mini{Representation tradeoff}
\begin{itemize}
\item byte: finite vocab, long seqs
\item word: short seqs, brittle OOV/morphology
\item subword: practical middle ground
\end{itemize}

\mini{Idealized objective}
\defs{\(\Sigma=\) vocab,\; \(f=\) tokenizer/segmentation fn,\; \(V=\) target vocab size.}
\[
\min_{\Sigma,f}\sum_i |f(x^{(i)})|
\quad
\text{s.t. } |\Sigma|=V
\]
\intuition{Compress the corpus into few tokens while respecting a vocab budget.}

\mini{BPE}
\begin{itemize}
\item start from bytes/chars; repeatedly merge most frequent adjacent pair
\item store merge list; apply merges greedily at test time
\end{itemize}

\mini{Other families}
\begin{itemize}
\item WordPiece: likelihood / PMI flavored
\item Unigram LM: start big, prune
\item SentencePiece: BPE / Unigram without whitespace assumptions
\end{itemize}

\mini{Keep straight}
\begin{itemize}
\item pretokenization matters; compression \(\neq\) downstream quality
\item no universal winner; special chat tokens stay atomic
\end{itemize}

\end{multicols}

\end{document}
